\begin{center}
\textsc{Gauss} an \textsc{Olbers}.\quad Göttingen, 2.~April 1826.
\end{center}

\par\noindent\hspace*{\parindent}\mbox{.\ .\ .\ .\ .\ .}\quad Was meine Messungen betrifft, kann ich \textls[120]{für mich allein} wenig entscheiden.
Mein \textls[120]{Auftrag} ist, was den trigonometrischen Teil betrifft, im Grunde abgeschlossen, und ich dachte, sofern \textsc{Schumacher} mich \textls[120]{ausreichend} unterstützen will, im Spätsommer die Beobachtungen mit dem Zenitsektor vorzunehmen.
Mein schwankender Gesundheitszustand schreckt mich davor ab, weitergehende trigonometrische Messungen vorzuschlagen. Unterdessen habe ich vor Kurzem, \textls[120]{unter uns gesagt}, in einem Schreiben an \textsc{Münster} erklärt, dass ich bereit bin, auch künftig meine Kräfte dafür einzusetzen, falls solche Messungen verlangt werden sollten.

\par\noindent\hspace*{\parindent}Meine theoretischen Arbeiten haben bei ihrem so \textls[120]{sehr} großen Umfang leider noch viele Lücken. Am einfachsten wäre mir geholfen, wenn ich mir
\par\clearpage
\noindent\null\hfill{\small\textbf{ZUR HANNOVERSCHEN TRIANGULATION.}}\hfill\makebox[0pt][r]{\fontsize{13}{16}\selectfont \textbf{377}}\par\medskip
\noindent erlaubte, mit der Veröffentlichung meiner Messungen zwar alle meine Rechenverfahren zu verbinden, ihre \textls[120]{Herleitung} aus den tiefer liegenden Gründen aber für ein ganz eigenständiges Werk in glücklicheren künftigen Zeiten aufzubewahren.
Dann gäbe es nirgends ein Hindernis.
Vorerst werde ich die präzise Ausgleichung meiner 32 Punkte vornehmen, die 51 Dreiecke und 146 Richtungen ergeben.
Die Höhenausgleichung (eine sehr viel leichtere Aufgabe) habe ich in diesen Tagen abgeschlossen.\quad\mbox{.\ .\ .\ .\ .\ .}
